\documentclass[addpoints]{exam}
\usepackage{amsmath}
\usepackage{amssymb}
\usepackage{color}
\usepackage{booktabs}

% \printanswers

\newcommand{\event}{Modelling with Functions}
\newcommand{\tournament}{}
\def\type{0}

\setlength\fillinlinelength{\textwidth}
\CorrectChoiceEmphasis{\color{red}\bfseries}
\SolutionEmphasis{\color{red}\bfseries}

\setlength\answerclearance{2pt}
\pagestyle{headandfoot}
\runningheadrule
\runningheader{\event}{\tournament}{\ifnum\type=1 Class Set Test \else Full Name:\hspace{1cm} \fi}
\runningfootrule
\runningfooter{}{Page \thepage\ of \numpages}{}

\begin{document}
\begin{center}
    {\huge Grade 12 Foundations for College Math \\ \event
    \ifprintanswers \space Key
    \else \space \ifnum\type<2 Test \fi \ifnum\type=1 \textbf{(CLASS SET)} \fi
          \ifnum\type=2 Answer Sheet \fi
    \fi \\ \vspace{3mm}\tournament\vspace{1mm}}
\end{center}

\vspace{.25\textheight}

\begin{center}
    First Name: \ifprintanswers
    \underline{\hspace{3.5cm}}\textbf{KEY}\underline{\hspace{5.5cm}}\\\ \vspace{5mm}
    \else \underline{\hspace{10cm}}\\ \vspace{5mm} \fi
    Last Name: \underline{\hspace{10cm}}\\ \vspace{5mm}
\end{center}

\noindent\textbf{\underline{Directions:}}
\begin{itemize}
    \ifnum\type=1 \item \textbf{This test is a class set.} Do not write on this paper. \fi
    \item Show all work for full marks. Round dollar amounts to the nearest cent.
    \item The test is designed to be completed in 75 minutes.
\end{itemize}
\begin{center}
\textit{For grading use only}\\
\multirowgradetable{1}[pages]
\end{center}

\newpage
\begin{questions}

\section*{Part A --- Multiple Choice (15 marks)}
\vspace{5mm}

% Q1
\question[1]
A scatter plot of data shows a curved shape that rises steeply, with no maximum. The best model is:
\begin{choices}
    \choice Linear
    \choice Quadratic
    \correctchoice Exponential
    \choice Sinusoidal
\end{choices}
\vspace{3mm}

% Q2
\question[1]
The correlation coefficient $r$ for a strong positive linear relationship is closest to:
\begin{choices}
    \choice $-0.95$
    \choice $0.10$
    \choice $0.50$
    \correctchoice $0.97$
\end{choices}
\vspace{3mm}

% Q3
\question[1]
The vertex form of a quadratic function is:
\begin{choices}
    \correctchoice $f(x) = a(x-h)^2 + k$
    \choice $f(x) = ax^2 + bx + c$
    \choice $f(x) = a \cdot b^x$
    \choice $f(x) = a\sin(k(x-d)) + c$
\end{choices}
\vspace{3mm}

% Q4
\question[1]
An exponential growth model has the form $A = A_0 b^t$ where $b > 1$. If the doubling time is 5 years, then $b$ satisfies:
\begin{choices}
    \choice $b^{10} = 2$
    \correctchoice $b^5 = 2$
    \choice $b = 2^5$
    \choice $5b = 2$
\end{choices}
\vspace{3mm}

% Q5
\question[1]
A sinusoidal function $f(x) = a\sin(k(x-d)) + c$ has amplitude 3 and vertical shift up 5. Its maximum value is:
\begin{choices}
    \choice 3
    \choice 5
    \correctchoice 8
    \choice 2
\end{choices}
\vspace{3mm}

% Q6
\question[1]
Which statement about extrapolation is correct?
\begin{choices}
    \choice Extrapolation is always reliable for exponential models
    \correctchoice Extrapolation can be unreliable because the trend may not continue outside the data range
    \choice Extrapolation only applies to linear models
    \choice Extrapolation is the same as interpolation
\end{choices}
\vspace{3mm}

% Q7
\question[1]
The period of $f(x) = \sin(3x)$ is:
\begin{choices}
    \choice $3\pi$
    \choice $2\pi$
    \correctchoice $\dfrac{2\pi}{3}$
    \choice $\dfrac{\pi}{3}$
\end{choices}
\vspace{3mm}

% Q8
\question[1]
A quadratic model $h(t) = -4.9t^2 + 20t + 1$ represents the height (m) of a projectile. The maximum height occurs at:
\begin{choices}
    \choice $t = 20$
    \correctchoice $t \approx 2.04\ \text{s}$
    \choice $t = 4.9\ \text{s}$
    \choice $t = 1\ \text{s}$
\end{choices}
\vspace{3mm}

% Q9
\question[1]
Data shows a population that triples every 10 years. Starting from 1{,}000, the population after 30 years is:
\begin{choices}
    \choice 9{,}000
    \choice 27{,}000
    \correctchoice 27{,}000
    \choice 3{,}000
\end{choices}
\vspace{3mm}

% Q10
\question[1]
The amplitude of a sinusoidal function is calculated as:
\begin{choices}
    \choice max $-$ min
    \choice $\dfrac{\text{max} + \text{min}}{2}$
    \correctchoice $\dfrac{\text{max} - \text{min}}{2}$
    \choice max $+$ min
\end{choices}
\vspace{3mm}

% Q11
\question[1]
A linear regression gives $y = 2.5x + 10$ with $r = 0.88$. Which statement is correct?
\begin{choices}
    \choice The model explains 88\% of the variance in $y$
    \correctchoice There is a strong positive linear relationship between the variables
    \choice The slope is 10
    \choice The data must be quadratic
\end{choices}
\vspace{3mm}

% Q12
\question[1]
An exponential decay function has the form $A = A_0(0.8)^t$. The percent decrease per time period is:
\begin{choices}
    \choice 80\%
    \correctchoice 20\%
    \choice 8\%
    \choice 0.8\%
\end{choices}
\vspace{3mm}

% Q13
\question[1]
The vertical shift $c$ in a sinusoidal model represents:
\begin{choices}
    \choice The amplitude
    \correctchoice The midline (average value) of the function
    \choice The period
    \choice The phase shift
\end{choices}
\vspace{3mm}

% Q14
\question[1]
Given the quadratic $f(x) = -2(x-3)^2 + 18$, the $x$-intercepts are found by setting $f(x) = 0$. They are:
\begin{choices}
    \choice $x = 3$ only
    \choice $x = 0$ and $x = 6$
    \correctchoice $x = 0$ and $x = 6$
    \choice $x = 3 \pm 3$
\end{choices}
\vspace{3mm}

% Q15
\question[1]
A model predicts that a city's population was 50{,}000 in 2000 and grows at 2.5\% per year. The equation is:
\begin{choices}
    \choice $P = 50{,}000 + 2.5t$
    \choice $P = 50{,}000(2.5)^t$
    \correctchoice $P = 50{,}000(1.025)^t$
    \choice $P = 50{,}000(0.975)^t$
\end{choices}
\vspace{3mm}

\section*{Part B --- Short Answer (33 marks)}

\vspace{5mm}

%% Q16 --- Choosing a model
\question[5]
\textbf{Choosing a Model.}
The table below shows the number of bacteria $N$ in a culture over time $t$ (hours).

\begin{center}
\begin{tabular}{@{}cc@{}}
\toprule
$t$ (h) & $N$ (bacteria) \\
\midrule
0 & 100 \\
1 & 200 \\
2 & 400 \\
3 & 800 \\
4 & 1600 \\
\bottomrule
\end{tabular}
\end{center}

\begin{parts}
\part[2] Determine which type of model (linear, quadratic, or exponential) best fits this data. Justify your choice with at least two observations from the table.
\part[3] Write an equation for the model. Use it to predict the number of bacteria at $t = 6$ hours.
\end{parts}

\begin{solution}
\textbf{(a)} The data is best modelled by an \textbf{exponential} function. Justification:
\begin{itemize}
    \item The values are \textit{doubling} each hour (constant ratio of 2), not increasing by a constant amount (linear) or by a growing fixed amount (quadratic).
    \item A second difference test for a quadratic would not be constant; the first ratio (200/100, 400/200, 800/400, 1600/800) is always 2.
\end{itemize}

\textbf{(b)} The initial value is $N_0 = 100$ and the base is $b = 2$:
\[
N(t) = 100 \cdot 2^t
\]
At $t = 6$:
\[
N(6) = 100 \cdot 2^6 = 100 \times 64 = 6{,}400\ \text{bacteria}
\]
\end{solution}

\vspace{5mm}

%% Q17 --- Linear regression
\question[6]
\textbf{Linear Regression.}
A study records the number of hours studied ($x$) and the exam score ($y$) for 5 students:

\begin{center}
\begin{tabular}{@{}cc@{}}
\toprule
Hours studied ($x$) & Exam score ($y$) \\
\midrule
1 & 55 \\
2 & 62 \\
3 & 70 \\
4 & 75 \\
5 & 83 \\
\bottomrule
\end{tabular}
\end{center}

The line of best fit is found to be $y = 7x + 47$ with $r = 0.996$.

\begin{parts}
\part[2] Interpret the slope and $y$-intercept in context.
\part[2] Predict the exam score for a student who studies 7 hours. Comment on the reliability of this prediction.
\part[2] A student scored 90. Estimate how many hours they studied. Is this interpolation or extrapolation?
\end{parts}

\begin{solution}
\textbf{(a)} \textit{Slope} $= 7$: For each additional hour studied, the exam score increases by approximately 7 marks.\\
\textit{$y$-intercept} $= 47$: A student who studies 0 hours is predicted to score 47 (though this extrapolation may not be meaningful in practice).

\textbf{(b)} At $x = 7$: $y = 7(7) + 47 = 49 + 47 = 96$.\\
This is \textit{extrapolation} (beyond $x = 5$, the range of data). Since $r = 0.996$ the linear trend is very strong, but trends can change outside the observed range; the prediction should be treated with caution.

\textbf{(c)} Set $y = 90$: $90 = 7x + 47 \implies 7x = 43 \implies x \approx 6.1\ \text{hours}$.\\
This is \textit{extrapolation} (slightly beyond the data range of 1--5 hours).
\end{solution}

\vspace{5mm}

%% Q18 --- Quadratic model (projectile)
\question[6]
\textbf{Quadratic Model --- Thrown Ball.}
A ball is thrown upward from a platform. Its height $h$ (in metres) at time $t$ (in seconds) is given in the table:

\begin{center}
\begin{tabular}{@{}cc@{}}
\toprule
$t$ (s) & $h$ (m) \\
\midrule
0 & 5.0 \\
1 & 20.1 \\
2 & 25.4 \\
3 & 20.9 \\
4 & 6.6 \\
\bottomrule
\end{tabular}
\end{center}

A quadratic regression gives $h(t) = -5t^2 + 20t + 5$.

\begin{parts}
\part[2] Identify the vertex and state the maximum height and when it occurs.
\part[2] Find when the ball hits the ground (i.e., when $h = 0$). Round to two decimal places.
\part[2] What does the $h$-intercept represent in context?
\end{parts}

\begin{solution}
\textbf{(a)} The function is $h(t) = -5t^2 + 20t + 5 = -5(t^2 - 4t) + 5 = -5(t-2)^2 + 20 + 5 = -5(t-2)^2 + 25$.\\
Vertex: $(2,\ 25)$. The \textbf{maximum height is 25 m}, occurring at $t = 2\ \text{s}$.

\textbf{(b)} Set $h(t) = 0$:
\[
-5t^2 + 20t + 5 = 0 \implies t^2 - 4t - 1 = 0 \implies t = \frac{4 \pm \sqrt{16+4}}{2} = \frac{4 \pm \sqrt{20}}{2} = 2 \pm \sqrt{5}
\]
$t = 2 + \sqrt{5} \approx 2 + 2.236 = 4.24\ \text{s}$ (taking the positive root).\\
The ball hits the ground at approximately $t \approx 4.24\ \text{s}$.

\textbf{(c)} The $h$-intercept $(0,\ 5)$ represents the \textit{initial height} of the ball at $t = 0$, meaning the ball was released from a height of 5 m above the ground (on a platform).
\end{solution}

\vspace{5mm}

%% Q19 --- Exponential model / population
\question[8]
\textbf{Exponential Model --- World Population.}
World population data (approximate) is given below:

\begin{center}
\begin{tabular}{@{}cc@{}}
\toprule
Year & Population (billions) \\
\midrule
1960 & 3.0 \\
1975 & 4.1 \\
1990 & 5.3 \\
2005 & 6.5 \\
2020 & 7.8 \\
\bottomrule
\end{tabular}
\end{center}

Let $t$ be the number of years after 1960.

\begin{parts}
\part[2] Explain why an exponential model $P(t) = P_0 \cdot b^t$ is reasonable for this data. Use the data to estimate $b$ (the annual growth factor) by using two well-separated data points.
\part[2] Write a complete exponential equation for $P(t)$.
\part[2] Estimate the doubling time: find $t$ such that $P(t) = 6.0$ billion (double the 1960 value) using your equation.
\part[2] Use your model to predict the world population in 2040. Discuss one limitation of your prediction.
\end{parts}

\begin{solution}
\textbf{(a)} The ratios of consecutive values are roughly constant (e.g.\ $4.1/3.0 \approx 1.367$ over 15 years, $5.3/4.1 \approx 1.293$ over 15 years), suggesting approximately proportional growth --- consistent with exponential growth.

Using $t=0$ ($P=3.0$) and $t=60$ ($P=7.8$):
\[
7.8 = 3.0 \cdot b^{60} \implies b^{60} = 2.6 \implies b = 2.6^{1/60} = e^{\ln(2.6)/60} \approx e^{0.01596} \approx 1.01609
\]
Annual growth factor $b \approx 1.0161$ (about 1.61\% per year).

\textbf{(b)}
\[
P(t) = 3.0 \times (1.0161)^t \quad \text{(billions, } t = \text{years after 1960)}
\]

\textbf{(c)} Solve $3.0(1.0161)^t = 6.0$:
\[
(1.0161)^t = 2 \implies t = \frac{\ln 2}{\ln 1.0161} \approx \frac{0.6931}{0.01597} \approx 43.4\ \text{years}
\]
The population doubles (reaches 6 billion) approximately 43 years after 1960, i.e., around \textbf{2003}.

\textbf{(d)} $t = 2040 - 1960 = 80$:
\[
P(80) = 3.0 \times (1.0161)^{80} \approx 3.0 \times e^{80 \times 0.01596} \approx 3.0 \times e^{1.277} \approx 3.0 \times 3.586 \approx 10.8\ \text{billion}
\]
\textit{Limitation:} Growth rates can change due to factors such as falling birth rates, policy changes, or resource constraints. A constant growth-rate assumption may overestimate future populations.
\end{solution}

\vspace{5mm}

%% Q20 --- Sinusoidal model
\question[8]
\textbf{Sinusoidal Model --- Monthly Temperatures.}
The table shows the average monthly temperature (°C) in a Canadian city:

\begin{center}
\begin{tabular}{@{}lc@{}}
\toprule
Month & Avg Temp (°C) \\
\midrule
January   & $-15$ \\
February  & $-12$ \\
March     & $-3$ \\
April     & $7$ \\
May       & $14$ \\
June      & $19$ \\
July      & $22$ \\
August    & $20$ \\
September & $13$ \\
October   & $5$ \\
November  & $-5$ \\
December  & $-13$ \\
\bottomrule
\end{tabular}
\end{center}

\begin{parts}
\part[2] Find the amplitude $a$, the period, and the vertical shift (midline) $c$.
\part[2] Determine the value of $k$ in $T(m) = a\sin(k(m - d)) + c$ where $m$ is the month number ($m = 1$ for January).
\part[2] Estimate the phase shift $d$ and write the complete equation.
\part[2] Use your equation to predict the temperature in month $m = 10$ (October). Compare with the table value.
\end{parts}

\begin{solution}
\textbf{(a)}
Max temperature $= 22$°C (July); Min temperature $= -15$°C (January).
\[
a = \frac{22-(-15)}{2} = \frac{37}{2} = 18.5\ \text{°C}
\]
\[
c = \frac{22+(-15)}{2} = \frac{7}{2} = 3.5\ \text{°C}
\]
Period $= 12$ months (one year).

\textbf{(b)}
\[
k = \frac{2\pi}{12} = \frac{\pi}{6}
\]

\textbf{(c)} The maximum occurs at $m = 7$ (July). For a sine function the maximum occurs when the argument equals $\pi/2$:
\[
\frac{\pi}{6}(7 - d) = \frac{\pi}{2} \implies 7 - d = 3 \implies d = 4
\]
Complete equation:
\[
\boxed{T(m) = 18.5\sin\!\left(\frac{\pi}{6}(m - 4)\right) + 3.5}
\]

\textbf{(d)} At $m = 10$:
\[
T(10) = 18.5\sin\!\left(\frac{\pi}{6}(10-4)\right) + 3.5
       = 18.5\sin(\pi) + 3.5
       = 18.5(0) + 3.5 = 3.5\ \text{°C}
\]
The table gives 5°C; the model predicts 3.5°C, which is reasonably close. The small difference is due to the data not being a perfect sinusoid.
\end{solution}

\end{questions}
\end{document}
